% Источники: 07_Программные_требования/DOC-06_Программные_требования.md
% Генератор: docs/build_pdf.py
% Шаблон печатного представления. Содержимое подставляет docs/build_pdf.py.
\documentclass[12pt,a4paper]{article}
\usepackage{iftex}
\ifPDFTeX
  \usepackage[utf8]{inputenc}
  \usepackage[T2A]{fontenc}
  \usepackage{newtxtext,newtxmath}
\else
  \usepackage{fontspec}
  \IfFontExistsTF{Times New Roman}
    {\setmainfont{Times New Roman}}
    {\setmainfont{Libertinus Serif}}
  \IfFontExistsTF{Arial}{\setsansfont{Arial}}{\setsansfont{Latin Modern Sans}}
  \IfFontExistsTF{Menlo}{\setmonofont{Menlo}[Scale=MatchLowercase]}
    {\setmonofont{Latin Modern Mono}[Scale=MatchLowercase]}
\fi
\usepackage[russian]{babel}
\usepackage{geometry}
\geometry{left=25mm,right=15mm,top=20mm,bottom=20mm}
\usepackage{setspace}
\onehalfspacing
\usepackage{indentfirst}
\setlength{\parindent}{1.25cm}
\setlength{\parskip}{2pt}
\setlength{\emergencystretch}{3em}
\usepackage{graphicx,booktabs,longtable,array,tabularx,calc,etoolbox}
\usepackage{enumitem}
\setlist[itemize]{leftmargin=1.25cm,itemsep=2pt,topsep=4pt}
\setlist[enumerate]{leftmargin=1.25cm,itemsep=2pt,topsep=4pt}
\newcommand{\tightlist}{\setlength{\itemsep}{0pt}\setlength{\parskip}{0pt}}
\AtBeginEnvironment{longtable}{\small\singlespacing}
\setlength{\LTpre}{8pt}
\setlength{\LTpost}{8pt}
\setlength{\tabcolsep}{5pt}
\setlength{\LTleft}{0pt}
\setlength{\LTright}{0pt}
\usepackage{xurl}
\usepackage[unicode,colorlinks=true,linkcolor=black,urlcolor=blue]{hyperref}
\usepackage{bookmark}
\hypersetup{pdftitle={Сводные программные требования},pdfauthor={Анохин Г.А.},pdfsubject={Криптотренажёр mephi-12}}
\urlstyle{same}
\pagestyle{plain}
\setcounter{secnumdepth}{3}
\setcounter{tocdepth}{2}

\begin{document}
\hypersetup{pageanchor=false}
\begin{titlepage}
  \thispagestyle{empty}
  \centering
  {\large НИЯУ МИФИ\par}
  \vspace*{3.5cm}
  {\LARGE\bfseries Сводные программные требования\par}
  \vspace{0.5cm}
  {\large DOC-06\par}
  \vspace{0.6cm}
  {\large Платформа для изучения и практики курса по криптографии\par}
  {\large Криптотренажёр mephi-12\par}
  \vspace{0.7cm}
  {\large Версия 0.5\par}
  {\large Редакция от 07.11.2026\par}
  \vspace{1cm}
  \begin{tabularx}{\textwidth}{@{}p{3.4cm}Xp{4cm}@{}}
    \toprule
    Роль & ФИО & Подпись / Дата \\
    \midrule
    Разработал & Анохин Г.А. & \\
    \addlinespace[6pt]
    Согласовано & & \\
    \addlinespace[6pt]
    Утверждено & & \\
    \bottomrule
  \end{tabularx}
  \vfill
  {\large Москва, 2026\par}
\end{titlepage}
\hypersetup{pageanchor=true}

\section*{Лист регистрации изменений}
\addcontentsline{toc}{section}{Лист регистрации изменений}
\begingroup\small\singlespacing
\noindent
\begin{tabularx}{\textwidth}{@{}p{1.2cm}p{2cm}p{2.5cm}Xp{1.8cm}@{}}
\toprule
Версия & Дата & Автор & Содержание изменения & Утвердил \\
\midrule
0.5 & 07.11.2026 & Анохин Г.А. & Уточнены правила и контракты; расширено описание сценариев; подготовлена печатная редакция. & \\
\bottomrule
\end{tabularx}
\endgroup
\clearpage
\begingroup\small\singlespacing
\tableofcontents
\endgroup
\clearpage

\hypertarget{doc-06-сводные-программные-требования}{}

{\def\LTcaptype{none} % do not increment counter
\begin{longtable}[]{@{}
  >{\raggedright\arraybackslash}p{(\linewidth - 2\tabcolsep) * \real{0.3200}}
  >{\raggedright\arraybackslash}p{(\linewidth - 2\tabcolsep) * \real{0.6800}}@{}}
\toprule\noalign{}
\begin{minipage}[b]{\linewidth}\raggedright
Поле
\end{minipage} & \begin{minipage}[b]{\linewidth}\raggedright
Значение
\end{minipage} \\
\midrule\noalign{}
\endhead
\bottomrule\noalign{}
\endlastfoot
Версия & 0.5 \\
Дата & 07.11.2026 \\
Редактор & Анохин Г.А. \\
Основание & Рабочая редакция DOC-05 \\
\end{longtable}
}

\section{Назначение и границы}\label{1-ux43dux430ux437ux43dux430ux447ux435ux43dux438ux435-ux438-ux433ux440ux430ux43dux438ux446ux44b}

Документ распределяет общие системные обязательства между логическими программными компонентами и определяет условия их совместной работы. Подсистемные требования уточняют эти обязательства; в текущую редакцию включён \href{DOC-06.2_\%D0\%A1\%D0\%B5\%D1\%81\%D1\%81\%D0\%B8\%D0\%B8_\%D0\%B7\%D0\%B0\%D0\%B4\%D0\%B0\%D1\%87_\%D0\%B8_\%D0\%BF\%D0\%BB\%D0\%B0\%D0\%BD\%D0\%B8\%D1\%80\%D0\%BE\%D0\%B2\%D1\%89\%D0\%B8\%D0\%BA.pdf\#nameddest=doc-062-\%D1\%81\%D0\%B5\%D1\%81\%D1\%81\%D0\%B8\%D0\%B8-\%D0\%B7\%D0\%B0\%D0\%B4\%D0\%B0\%D1\%87-\%D0\%B8-\%D0\%BF\%D0\%BB\%D0\%B0\%D0\%BD\%D0\%B8\%D1\%80\%D0\%BE\%D0\%B2\%D1\%89\%D0\%B8\%D0\%BA}{DOC-06.2 по сессиям}.

Компонент здесь означает область ответственности ПО. Такое разделение не требует отдельных сервисов или применения единственного архитектурного стиля. Математические алгоритмы, модель хранения и полный справочник HTTP-операций не дублируются в общей части.

\section{Ответственность и связи}\label{2-ux43eux442ux432ux435ux442ux441ux442ux432ux435ux43dux43dux43eux441ux442ux44c-ux438-ux441ux432ux44fux437ux438}

{\def\LTcaptype{none} % do not increment counter
\begin{longtable}[]{@{}
  >{\raggedright\arraybackslash}p{(\linewidth - 4\tabcolsep) * \real{0.2468}}
  >{\raggedright\arraybackslash}p{(\linewidth - 4\tabcolsep) * \real{0.4774}}
  >{\raggedright\arraybackslash}p{(\linewidth - 4\tabcolsep) * \real{0.2759}}@{}}
\toprule\noalign{}
\begin{minipage}[b]{\linewidth}\raggedright
Компонент
\end{minipage} & \begin{minipage}[b]{\linewidth}\raggedright
Вход и результат
\end{minipage} & \begin{minipage}[b]{\linewidth}\raggedright
Ответственный за предметную часть
\end{minipage} \\
\midrule\noalign{}
\endhead
\bottomrule\noalign{}
\endlastfoot
Проверка доступа & Учётные данные запроса; подтверждённый пользователь и полномочия либо отказ & Кокорев \\
Конфигурации & Пользователь и его группа; доступный состав набора либо отказ & Самбьюи \\
Управление сессиями & Запрос создания или ответа; сессия и предметный результат либо категория отказа & Анохин \\
Генераторы и проверяющие & Тип и настройки; условие. Сохранённое условие и ответ; предметный вердикт либо ошибка проверки & Владельцы криптографических доменов \\
Контроль срока & Сохранённые времена и сведения о принятых ответах; завершение просроченных задач & Анохин \\
Хранение & Данные и изменения состояния; согласованное сохранение и чтение & Анохин для правил сессий, Самбьюи для общей модели и инфраструктуры \\
\end{longtable}
}

При создании набора управление сессиями получает пользователя и доступную конфигурацию, вызывает генераторы и сохраняет состав. При отправке ответа оно проверяет допуск, обращается к профильному проверяющему и сохраняет результат. Контроль срока использует те же данные времени и завершения; отдельного правила подсчёта итога у него нет.

\subsection{Последовательность создания набора}\label{ux43fux43eux441ux43bux435ux434ux43eux432ux430ux442ux435ux43bux44cux43dux43eux441ux442ux44c-ux441ux43eux437ux434ux430ux43dux438ux44f-ux43dux430ux431ux43eux440ux430}

Компонент приёма операции передаёт управлению сессиями подтверждённого пользователя и имя конфигурации. Затем определяется группа пользователя и запрашивается доступный состав набора. Эти действия не заменяются доверием к переданному клиентом идентификатору группы.

По составу вызываются соответствующие генераторы. Полученные условия, типы и принадлежность задач образуют данные создаваемой сессии. Подтверждение выдаётся после сохранения полного набора. При отказе генерации или сохранения частичный набор не должен оставаться доступной активной сессией.

Ограничение одной активной сессии охватывает всю операцию создания. Отдельная предварительная проверка без согласования последующего сохранения недостаточна: два обращения могли бы одновременно увидеть отсутствие сессии. Конкретный механизм защиты выбирается при проектировании хранения, но программное обязательство уже определено.

\subsection{Последовательность обработки ответа}\label{ux43fux43eux441ux43bux435ux434ux43eux432ux430ux442ux435ux43bux44cux43dux43eux441ux442ux44c-ux43eux431ux440ux430ux431ux43eux442ux43aux438-ux43eux442ux432ux435ux442ux430}

Приём запроса фиксирует серверное время его поступления. Управление сессиями устанавливает задачу, проверяет принадлежность, состояние и срок, затем проверяет формат по схеме сохранённого типа. Предметный обработчик получает только допустимый ответ с условием выбранной задачи.

После получения вердикта управление сессиями согласованно сохраняет результат задачи и вычисляет итог набора. Контроль срока использует те же правила завершения и сведения о своевременном запросе. Если ответ получен до срока, момент окончания вычисления не отменяет его своевременность.

Общие требования к обмену определяют смысл данных на этих переходах. Поля и HTTP-коды раскрываются в интерфейсных спецификациях; математические критерии правильности --- в предметных подсистемах. Такое распределение позволяет уточнять внутреннее устройство части, сохраняя её обязательства перед остальными участниками.

\section{Общие программные требования}\label{3-ux43eux431ux449ux438ux435-ux43fux440ux43eux433ux440ux430ux43cux43cux43dux44bux435-ux442ux440ux435ux431ux43eux432ux430ux43dux438ux44f}

Все требования ниже обязательны. Ссылки SR ведут к \href{../04_\%D0\%A1\%D0\%B8\%D1\%81\%D1\%82\%D0\%B5\%D0\%BC\%D0\%BD\%D1\%8B\%D0\%B5_\%D1\%82\%D1\%80\%D0\%B5\%D0\%B1\%D0\%BE\%D0\%B2\%D0\%B0\%D0\%BD\%D0\%B8\%D1\%8F/DOC-05_\%D0\%A1\%D0\%B8\%D1\%81\%D1\%82\%D0\%B5\%D0\%BC\%D0\%BD\%D1\%8B\%D0\%B5_\%D1\%82\%D1\%80\%D0\%B5\%D0\%B1\%D0\%BE\%D0\%B2\%D0\%B0\%D0\%BD\%D0\%B8\%D1\%8F.pdf\#nameddest=doc-05-\%D1\%82\%D1\%80\%D0\%B5\%D0\%B1\%D0\%BE\%D0\%B2\%D0\%B0\%D0\%BD\%D0\%B8\%D1\%8F-\%D0\%BA-\%D0\%B8\%D0\%BD\%D1\%82\%D0\%B5\%D0\%B3\%D1\%80\%D0\%B0\%D1\%86\%D0\%B8\%D0\%B8}{системному разделу интеграции}; проверки являются описательными.

\subsection{SWR-SYS-001. Передача личности пользователя}\label{swr-sys-001-ux43fux435ux440ux435ux434ux430ux447ux430-ux43bux438ux447ux43dux43eux441ux442ux438-ux43fux43eux43bux44cux437ux43eux432ux430ux442ux435ux43bux44f}

Компоненты пользовательских операций должны получать личность пользователя из результата проверки доступа. Параметры тела и строки запроса не должны заменять подтверждённого владельца операции.

Источник: SR-SYS-001. Проверка: подстановка чужого идентификатора не изменяет пользователя, от имени которого выполняется операция.

\subsection{SWR-SYS-002. Общие идентификаторы}\label{swr-sys-002-ux43eux431ux449ux438ux435-ux438ux434ux435ux43dux442ux438ux444ux438ux43aux430ux442ux43eux440ux44b}

Компоненты должны сохранять и передавать идентификаторы пользователя, сессии и задачи без повторного назначения при обмене между ними.

Источник: SR-SYS-002. Проверка: задача в выдаче, проверке и хранилище имеет одну принадлежность и один идентификатор.

\subsection{SWR-SYS-003. Состав созданной сессии}\label{swr-sys-003-ux441ux43eux441ux442ux430ux432-ux441ux43eux437ux434ux430ux43dux43dux43eux439-ux441ux435ux441ux441ux438ux438}

Управление сессиями должно сохранять состав и условия выданных задач по результату выбора доступной конфигурации; последующая проверка должна использовать эти сохранённые данные.

Источники: SR-SYS-003, SR-SYS-004. Проверка: изменение конфигурации после создания не меняет условия уже выданной задачи.

\subsection{SWR-SYS-004. Контракт предметного обработчика}\label{swr-sys-004-ux43aux43eux43dux442ux440ux430ux43aux442-ux43fux440ux435ux434ux43cux435ux442ux43dux43eux433ux43e-ux43eux431ux440ux430ux431ux43eux442ux447ux438ux43aux430}

Предметный обработчик должен различать генерацию условия и проверку ответа, а результат проверки --- завершённый вердикт, ошибку формата или технический сбой без вердикта.

Источники: SR-SYS-004, SR-SYS-009. Проверка: обработчик вызывается по сохранённому типу; сбой не возвращается как неверное решение. Форматы своих данных подтверждает владелец домена.

\subsection{SWR-SYS-005. Согласованное сохранение}\label{swr-sys-005-ux441ux43eux433ux43bux430ux441ux43eux432ux430ux43dux43dux43eux435-ux441ux43eux445ux440ux430ux43dux435ux43dux438ux435}

Изменение состояния задачи и вычисление соответствующего состояния сессии должны сохраняться как одна согласованная операция до выдачи подтверждённого результата.

Источник: SR-SYS-005. Проверка: сбой сохранения не оставляет подтверждённый пользователю вердикт с противоречащим ему состоянием сессии.

\subsection{SWR-SYS-006. Общий серверный срок}\label{swr-sys-006-ux43eux431ux449ux438ux439-ux441ux435ux440ux432ux435ux440ux43dux44bux439-ux441ux440ux43eux43a}

Обработка ответов и контроль просрочки должны использовать одно сохранённое время создания сессии, один срок и зафиксированное сервером время получения ответа.

Источники: SR-SYS-006, SR-SYS-007. Проверка: перезапуск не продлевает срок, а задержка планировщика не меняет допустимость ответа.

\subsection{SWR-SYS-007. Категория результата операции}\label{swr-sys-007-ux43aux430ux442ux435ux433ux43eux440ux438ux44f-ux440ux435ux437ux443ux43bux44cux442ux430ux442ux430-ux43eux43fux435ux440ux430ux446ux438ux438}

Компонент обработки запроса должен сохранять различие между предметным вердиктом и отказами при передаче результата внешнему интерфейсу.

Источники: SR-SYS-008, SR-SYS-009. Проверка: неверное решение, нарушение доступа, ошибка формата, просрочка и технический сбой распознаются как разные категории; соответствие HTTP-кодам задано контрактом сессий и сохраняется в общем справочнике DOC-08.

\subsection{SWR-SYS-008. Сохранность завершённых данных}\label{swr-sys-008-ux441ux43eux445ux440ux430ux43dux43dux43eux441ux442ux44c-ux437ux430ux432ux435ux440ux448ux451ux43dux43dux44bux445-ux434ux430ux43dux43dux44bux445}

Хранение должно обеспечивать восстановление идентификаторов, принадлежности, условий, времени создания и завершённых результатов сессий после перезапуска приложения.

Источник: SR-SYS-007. Проверка: после восстановления выдаётся прежний итог и применяется прежний срок. Для принятого незавершённого ответа используются сохранённая операция и прежний срок обработки по DOC-06.2.

\section{Состав подсистемных документов}\label{4-ux441ux43eux441ux442ux430ux432-ux43fux43eux434ux441ux438ux441ux442ux435ux43cux43dux44bux445-ux434ux43eux43aux443ux43cux435ux43dux442ux43eux432}

{\def\LTcaptype{none} % do not increment counter
\begin{longtable}[]{@{}
  >{\raggedright\arraybackslash}p{(\linewidth - 4\tabcolsep) * \real{0.3276}}
  >{\raggedright\arraybackslash}p{(\linewidth - 4\tabcolsep) * \real{0.3276}}
  >{\raggedright\arraybackslash}p{(\linewidth - 4\tabcolsep) * \real{0.3447}}@{}}
\toprule\noalign{}
\begin{minipage}[b]{\linewidth}\raggedright
Документ
\end{minipage} & \begin{minipage}[b]{\linewidth}\raggedright
Владелец
\end{minipage} & \begin{minipage}[b]{\linewidth}\raggedright
Предметная область
\end{minipage} \\
\midrule\noalign{}
\endhead
\bottomrule\noalign{}
\endlastfoot
DOC-06.1 & Кокорев & Аутентификация и пользователи \\
\href{DOC-06.2_\%D0\%A1\%D0\%B5\%D1\%81\%D1\%81\%D0\%B8\%D0\%B8_\%D0\%B7\%D0\%B0\%D0\%B4\%D0\%B0\%D1\%87_\%D0\%B8_\%D0\%BF\%D0\%BB\%D0\%B0\%D0\%BD\%D0\%B8\%D1\%80\%D0\%BE\%D0\%B2\%D1\%89\%D0\%B8\%D0\%BA.pdf\#nameddest=doc-062-\%D1\%81\%D0\%B5\%D1\%81\%D1\%81\%D0\%B8\%D0\%B8-\%D0\%B7\%D0\%B0\%D0\%B4\%D0\%B0\%D1\%87-\%D0\%B8-\%D0\%BF\%D0\%BB\%D0\%B0\%D0\%BD\%D0\%B8\%D1\%80\%D0\%BE\%D0\%B2\%D1\%89\%D0\%B8\%D0\%BA}{DOC-06.2} & Анохин & Сессии задач и планировщик \\
DOC-06.3 & Георгиев & Ранцевые задачи \\
DOC-06.4 & Минасов & Учебный бэкдор-сценарий \\
DOC-06.5 & Портнов & Эль-Гамаль \\
DOC-06.6 & Самбьюи & Конфигурации и модерация \\
\end{longtable}
}

\subsection{Обязательства на границах подсистем}\label{ux43eux431ux44fux437ux430ux442ux435ux43bux44cux441ux442ux432ux430-ux43dux430-ux433ux440ux430ux43dux438ux446ux430ux445-ux43fux43eux434ux441ux438ux441ux442ux435ux43c}

Безопасность задаёт подтверждённого пользователя и условия доступа; сессии используют этот результат. Конфигурации предоставляют доступный состав, по которому создаётся набор. Предметные обработчики различают вердикт и невозможность проверки; общий интерфейс сохраняет это различие при сообщении результата.

Задача, итог сессии и запись обработанного ответа сохраняются в одной транзакции. Незавершённая запись служит основанием восстановления. Время получения ответа, общий срок и предел обработки сохраняются при повторном запуске; вычислительная задержка не превращает своевременный ответ в поздний.

\section{Связи и проверка редакции}\label{5-ux441ux432ux44fux437ux438-ux438-ux43fux440ux43eux432ux435ux440ux43aux430-ux440ux435ux434ux430ux43aux446ux438ux438}

Требования SR-SYS-001--SR-SYS-009 раскрыты восемью SWR-SYS выше. Для каждого указаны источник и ожидаемая проверка. Детализация SR-SES находится в DOC-06.2. Описательные кейсы DOC-10 связываются с этими требованиями в матрице DOC-09.

Проверка редакции охватывает единый контракт ошибок, точность целых чисел, доступ по владельцу, согласованность сохранения и восстановление принятого ответа. \href{../04_\%D0\%A1\%D0\%B8\%D1\%81\%D1\%82\%D0\%B5\%D0\%BC\%D0\%BD\%D1\%8B\%D0\%B5_\%D1\%82\%D1\%80\%D0\%B5\%D0\%B1\%D0\%BE\%D0\%B2\%D0\%B0\%D0\%BD\%D0\%B8\%D1\%8F/DOC-05_\%D0\%A1\%D0\%B8\%D1\%81\%D1\%82\%D0\%B5\%D0\%BC\%D0\%BD\%D1\%8B\%D0\%B5_\%D1\%82\%D1\%80\%D0\%B5\%D0\%B1\%D0\%BE\%D0\%B2\%D0\%B0\%D0\%BD\%D0\%B8\%D1\%8F.pdf\#nameddest=5-\%D1\%80\%D0\%B0\%D1\%81\%D0\%BF\%D1\%80\%D0\%B5\%D0\%B4\%D0\%B5\%D0\%BB\%D0\%B5\%D0\%BD\%D0\%B8\%D0\%B5-\%D0\%BE\%D1\%82\%D0\%B2\%D0\%B5\%D1\%82\%D1\%81\%D1\%82\%D0\%B2\%D0\%B5\%D0\%BD\%D0\%BD\%D0\%BE\%D1\%81\%D1\%82\%D0\%B8}{DOC-05} задаёт распределение ответственности.

История: 05.11.2026 --- подготовлена версия 0.5: ответственность частей, общие программные обязательства, связи с системными требованиями и состав подсистем. 07.11.2026: расширено описание взаимодействий и определено восстановление; подготовлены LaTeX и PDF.


\end{document}
